\documentclass[a4paper,12pt]{article}
\usepackage{utils}

\begin{document}

\begin{titlepage}
    \centering
    \vspace*{2cm}

    {\scshape\Large Questions ouvertes \par}
    \vspace{2cm}

    \rule{\linewidth}{0.5mm} \\[0.4cm]
    {\huge\bfseries Compte rendu \par}
    \rule{\linewidth}{0.5mm} \\[1.5cm]

    \vspace{2cm}

    {\Large
    Lorenzo \textsc{Carrascosa} \par}

    \vfill
\end{titlepage}

\clearpage

\section*{Bilan de l'UE de questions ouvertes}

\subsection*{Points forts et pertinence du module}

Cette unité d'enseignement offre une véritable introduction à la recherche, comblant ainsi un manque d'un cursus dédié à la recherche.
Les sujets proposés sont intéressants, diversifiés et exploitent les outils acquis en première et deuxième année.

\subsection*{Les limites}

\begin{itemize}[label=\textbullet]
    \item \textbf{L'attribution des sujets} : Le nombre de sujets proposés devrait être supérieur au nombre de groupes, afin d'éviter que certains étudiants travaillent sur des sujets qu'ils n'ont pas choisis.
    
    \item \textbf{Le travail en groupe} : Il existe une forte disparité entre les étudiants intéressés par la recherche et ceux qui ne le sont pas, ce qui complique le travail en groupe.
    
    \item \textbf{Format d'évaluation} : Les modalités de restitution (journal de bord et oral de 10 minutes) sont insuffisantes pour rendre compte de la profondeur d'un travail de recherche et en garantir la qualité.
\end{itemize}

\subsection*{Suggestions d'amélioration}

\begin{itemize}[label=\textbullet]
    \item \textbf{Format individuel} : Permettre la rédaction d'un mémoire individuel, à l'image des TIPE en classe préparatoire (\href{https://licence-math.univ-lyon1.fr/doku.php?id=tipe:page}{à l'instar du modèle MAT2007L à lyon 1}). Ce format constituerait un atout décisif pour les dossiers d'intégration en Master, Magistère ou ENS.
\end{itemize}

\end{document}